\documentclass[11pt]{article}
\usepackage[margin=1in]{geometry}
\usepackage{amsmath, amssymb, bm}
\usepackage[T1]{fontenc}
\usepackage{lmodern}
\usepackage{hyperref}
\usepackage{enumitem}
\usepackage{booktabs}
\usepackage{array}
\usepackage{xurl}

\newcommand{\codepath}[1]{\path{#1}}

\title{Update Manual for \texttt{phase\_velocity.py}}
\author{RMFsolver package update notes}
\date{\today}

\begin{document}
\maketitle


\section{Purpose of This Manual}
This document records the recent steady-front updates implemented in
\codepath{RMFsolver/phase_velocity.py}. The recommended current production path
is the compact-coordinate boundary-value solver
\[
\texttt{solve\_steady\_front\_1d\_aQstar\_rescale\_bvp(...)}.
\]
This solver fixes the interface composition $a_{Q^\star}$, solves for the
compatible baryon flux $j_B$, and enforces the interface and downstream stable
tail conditions in a single BVP solve.

Older one-dimensional shooting and compact-IVP notes are kept in the \LaTeX{}
source for history, but they are wrapped in an \texttt{\textbackslash iffalse ...
\textbackslash fi} block. The rendered manual now focuses on the latest BVP
workflow and the helper functions needed to reuse it elsewhere.

The emphasis here is on:
\begin{itemize}[leftmargin=1.5em]
\item the current BVP method and why it replaced the intermediate shooting paths,
\item how the shared helper functions construct the nuclear, interface, and
      quark endpoint states,
\item how local hydro closure is performed without assuming constant quark-side
      velocity,
\item how branch selection is handled for the $\mu_K$-rich and $\mu_K$-poor
      solutions,
\item how to configure scan scripts so the BVP stays on the physical branch.
\end{itemize}

\section{High-Level Summary of the Update}
The steady-front portion of \texttt{phase\_velocity.py} now contains several
historical solver paths, but the current recommended path is
\texttt{solve\_steady\_front\_1d\_aQstar\_rescale\_bvp(...)}.
It treats $a_{Q^\star}$ as the scan variable and solves for $j_B$ together with
the functions $a(s)$ and $q(s)$ using \texttt{scipy.integrate.solve\_bvp}.

The older solver paths are still present in the Python file for comparison and
regression checks:
\begin{itemize}[leftmargin=1.5em]
\item \texttt{solve\_steady\_front\_2d(...)}: historical 2D shooting in
      $(j_B,a_{Q^\star})$;
\item \texttt{solve\_steady\_front\_1d\_aQstar(...)}: finite-domain scalar
      shooting in $j_B$ at fixed $a_{Q^\star}$;
\item \texttt{solve\_steady\_front\_1d\_aQstar\_rescaled(...)}:
      compact-coordinate IVP shooting in $j_B$.
\end{itemize}
Those intermediate 1D methods are archived in the source of this manual but are
not rendered, so the current document can serve as a clean reference for the BVP
implementation.

\section{Shared Physics Model Used by Both Solvers}
\subsection{Steady equations on the quark side}
On the quark side $(x > 0)$ the code solves the steady system
\begin{align}
j_B &= u(x)\, n_B(x), \\
\Pi &= h(x)\,u(x)^2 + P(x), \\
\partial_x \left(D\, \partial_x a - u a\right) &= R(a),
\end{align}
where the composition variable is
\begin{equation}
a = \frac{n_K - n_K^Q}{n_B^Q}.
\end{equation}

The flux variable is defined as
\begin{equation}
q(x) = D\, \partial_x a - u(x)\,a(x).
\end{equation}

The first-order system actually integrated by the code is
\begin{align}
\frac{da}{dx} &= \frac{q + u a}{D}, \\
\frac{dq}{dx} &= \gamma \left(a^3 + \eta a\right).
\end{align}

\subsection{Endpoint-state notation}
The code uses the following notation consistently:
\begin{itemize}[leftmargin=1.5em]
\item \textbf{N}: upstream nuclear state at $x=0^-$,
\item \textbf{$Q^\star$}: immediate quark-side interface state at $x=0^+$,
\item \textbf{Q}: far-right equilibrated quark state as $x \to \infty$.
\end{itemize}

\subsection{Frozen microphysics convention}
Both steady-front solvers use the same convention:
\begin{itemize}[leftmargin=1.5em]
\item $D$, $\eta$, and $\gamma$ are evaluated at $Q^\star$,
\item they are then kept fixed during the IVP integration.
\end{itemize}

\section{New and Modified Helper Functions}
\subsection{\texttt{nK\_QM(muB, muK, B\_one\_forth, T, ms=0, upB=5000, dmu=None)}}
This helper returns the kaon-density-like thermodynamic derivative
\[
n_K = \frac{\partial P_{QM}}{\partial \mu_K}.
\]

In the current implementation it is handled analytically from the quark flavor
densities:
\begin{equation}
n_K = \frac{n_d - n_s}{2}.
\end{equation}

The function calls the existing quark-density helper, obtains the physical
flavor densities $n_d$ and $n_s$, and returns their half-difference. This is
cleaner than the earlier finite-difference pressure derivative and avoids
introducing extra numerical noise into the steady-front state construction.

\subsection{\texttt{\_solve\_muB\_Q\_at\_muK0\_for\_given\_Pi(...)}} 
This helper solves for the far-right equilibrated quark endpoint $Q$ at
\[
\mu_K = 0
\]
for a given momentum flux $\Pi$ and baryon flux $j_B$.

Numerically it solves
\begin{equation}
\Pi_{QM}(\mu_B, 0, B^{1/4}, T, j_B) - \Pi = 0
\end{equation}
using \texttt{fsolve} with an initial guess of $1100$ MeV. After obtaining the
root, it checks that the corresponding quark baryon density is positive.

\subsection{\texttt{\_branch\_muK\_seed(a\_like)}}
This helper returns the positive-$\mu_K$ seed used by the current steady-front
quark-state solves:
\begin{equation}
\mu_K^{\rm seed} = \max(1,\ 20,\ 250 |a|),
\end{equation}
in MeV. This seed is used both at $Q^\star$ and during local continuation
along the profile. The current code does not implement a separate near-zero
$\mu_K$ branch mode in this helper.

\subsection{\texttt{\_solve\_interface\_Qstar\_from\_aQstar\_and\_Pi(...)}} 
This helper solves for the interface state $Q^\star$ by determining
$(\mu_B^{Q^\star}, \mu_K^{Q^\star})$ from the pair of equations
\begin{align}
\Pi_{QM}(\mu_B, \mu_K, B^{1/4}, T, j_B) &= \Pi, \\
\frac{n_K(\mu_B,\mu_K) - n_K^Q}{n_B^Q} &= a_{Q^\star}.
\end{align}

The implementation details are:
\begin{enumerate}[leftmargin=1.5em]
\item It first builds a list of initial guesses.
\item These guesses are all positive-$\mu_K$ seeds and include stronger
      positive guesses for larger $|a_{Q^\star}|$.
\item For each guess, it calls \texttt{scipy.optimize.root(..., method="hybr")}.
\item Successful roots are filtered to require positive density and
nonnegative $\mu_K$.
\item Near-duplicate roots are deduplicated.
\item The remaining candidates are ranked by closeness to the positive-$\mu_K$
      target seed and then by larger $\mu_K$, so that the solver stays on the
      intended positive-$\mu_K$ physical branch.
\end{enumerate}

This helper is central to maintaining branch consistency in the steady-front
machinery.

\subsection{\texttt{\_solve\_local\_quark\_state\_from\_a\_and\_Pi(...)}} 
This helper solves the local quark thermodynamic state at a given point in the
profile. Its unknowns are again $(\mu_B, \mu_K)$, but now the composition
constraint uses the local $a(x)$ instead of $a_{Q^\star}$.

It reuses
\texttt{\_solve\_interface\_Qstar\_from\_aQstar\_and\_Pi(...)} with
$a \leftrightarrow a_{Q^\star}$, then computes
\begin{equation}
n_B(x), \qquad u(x) = \frac{j_B}{n_B(x)}.
\end{equation}

In both the 2D and 1D solvers, the previous successful local state is passed as
the next initial guess. That gives the code a nearest-neighbor continuation
strategy along the profile and helps reduce accidental branch switching.

\subsection{\texttt{\_microphysics\_at\_Qstar(muB\_Qstar, T)}}
This helper computes the frozen microphysics coefficients using the current
implementation convention:
\begin{align}
\mu_{Q^\star} &= \mu_B^{Q^\star}/3, \\
q_D &= \sqrt{\frac{3 g_s^2 (\mu_{Q^\star})^2}{2\pi^2}}, \\
D &= \left[\frac{24 \alpha_s^2}{\pi}
\left(
h_{\rm const}\,\frac{T^{5/3}}{q_D^{2/3}}
+ \frac{\pi^3 T^2}{12 q_D}
\right)\right]^{-1}, \\
\eta &= \frac{9\pi^2 T^2}{(\mu_{Q^\star})^2}, \\
\tau &= 1.98\times 10^{12}\left(\frac{300}{\mu_{Q^\star}}\right)^5, \\
\gamma &= \frac{1}{\tau}.
\end{align}

These are returned in a dictionary and then treated as constants during the IVP.

\section{Behavior of the 2D Steady-Front Solver}
\subsection{Function}
\texttt{solve\_steady\_front\_2d(...)}

\subsection{Unknowns and residual}
The 2D solver still solves for both
\[
j_B, \qquad a_{Q^\star}
\]
using the outer shooting residual
\[
\bigl(a(x_{\max}), q(x_{\max})\bigr).
\]

Internally it solves in variables
\[
(\log j_B,\ a_{Q^\star})
\]
so that positivity of $j_B$ is automatic.

\subsection{Shared construction path}
For each shooting trial it:
\begin{enumerate}[leftmargin=1.5em]
\item builds the upstream nuclear state $N$,
\item solves the far-right equilibrated quark state $Q$ at $\mu_K=0$,
\item sets
      \[
      n_K^Q = 0
      \]
      exactly whenever $|m_s| \le 10^{-12}$,
\item computes
      \[
      a_N = \frac{n_B^N - n_K^Q}{n_B^Q},
      \]
\item solves the interface state $Q^\star$ using positive-$\mu_K$
      candidate-ranked root selection,
\item computes frozen $(D,\eta,\gamma,\tau)$ at $Q^\star$,
\item sets the initial condition
      \[
      q(0^+) = - a_N u_N,
      \]
\item integrates the IVP on $[0, x_{\max}]$.
\end{enumerate}

\subsection{Trivial-branch rejection}
The 2D solver also contains a guard against a near-trivial branch where:
\begin{itemize}[leftmargin=1.5em]
\item $a(0^+)$ is numerically tiny,
\item $q(0^+)$ is numerically tiny,
\item $Q^\star$ nearly collapses onto the equilibrated state $Q$.
\end{itemize}

% Archived intermediate 1D shooting / compact-IVP notes. Kept for history; not rendered.
\iffalse
\section{The New 1D Solver: \texttt{solve\_steady\_front\_1d\_aQstar}}
\subsection{Purpose}
The new function
\[
\texttt{solve\_steady\_front\_1d\_aQstar(T, nB\_N, B\_one\_forth, aQstar, ...)}
\]
treats $a_{Q^\star}$ as an input and solves only for $j_B$.

This is useful when the user wants to:
\begin{itemize}[leftmargin=1.5em]
\item hold the interface composition amplitude fixed,
\item scan $a_{Q^\star}$ as the preferred control parameter,
\item extract a one-dimensional solution curve
      \[
      j_B = j_B(a_{Q^\star}).
      \]
\end{itemize}

\subsection{Core idea}
Instead of demanding both
\[
a(x_{\max}) = 0, \qquad q(x_{\max}) = 0,
\]
the 1D solver uses the asymptotic decaying-tail condition at the outer
boundary.

Near equilibrium the system is linearized as
\begin{align}
\frac{da}{dx} &= \frac{q + u_Q a}{D}, \\
\frac{dq}{dx} &\approx \gamma \eta\, a.
\end{align}

For a decaying mode
\[
a(x) \sim e^{-\lambda x},
\]
the code uses
\begin{equation}
\lambda =
\frac{-u_Q + \sqrt{u_Q^2 + 4 D \gamma \eta}}{2 D},
\end{equation}
with
\[
u_Q = \frac{j_B}{n_B^Q}.
\]

The stable-tail relation is then
\begin{equation}
q \approx - (D \lambda + u_Q)\, a.
\end{equation}

Therefore the scalar residual used by the code is
\begin{equation}
F(j_B; a_{Q^\star}) =
q(x_{\max}) + \bigl(D \lambda + u_Q\bigr) a(x_{\max}).
\end{equation}

\subsection{Detailed algorithm}
For a fixed input $a_{Q^\star}$, the code does the following.

\subsubsection*{Step 1: choose a trial $j_B$}
The outer shooting variable is $j_B$, but internally the code uses
\[
z = \log j_B.
\]
This guarantees positivity and makes bounded scalar searches straightforward.

\subsubsection*{Step 2: build the trial state}
The nested helper \texttt{\_build\_trial\_state\_1d(...)} reuses the same state
construction path as the 2D solver:
\begin{enumerate}[leftmargin=1.5em]
\item construct the upstream nuclear state $N$,
\item compute
      \[
      \Pi = h_N u_N^2 + P_N, \qquad u_N = \frac{j_B}{n_B^N},
      \]
\item solve the far-right equilibrated quark state $Q$ at $\mu_K=0$,
\item set $n_K^Q = 0$ exactly if $m_s = 0$ numerically,
\item compute $a_N$,
\item solve the interface state $Q^\star$ using the branch-selected helper,
\item compute frozen microphysics at $Q^\star$,
\item integrate the IVP with initial data
      \[
      a(0^+) = a_{Q^\star}, \qquad q(0^+) = -a_N u_N.
      \]
\end{enumerate}

\subsubsection*{Step 3: integrate the IVP}
The code integrates
\begin{align}
\frac{da}{dx} &= \frac{q + u a}{D}, \\
\frac{dq}{dx} &= \gamma(a^3 + \eta a)
\end{align}
with \texttt{solve\_ivp(..., method="RK45")}.

At each RHS evaluation, the code solves the local quark state
$(\mu_B,\mu_K,n_B,u)$ consistently from the current $a(x)$ and the fixed
constants $(\Pi,j_B)$.

\subsubsection*{Step 4: compute the tail residual}
At the end of the integration, with
\[
a_L = a(x_{\max}), \qquad q_L = q(x_{\max}),
\]
the code computes
\begin{align}
u_Q &= \frac{j_B}{n_B^Q}, \\
\lambda &=
\frac{-u_Q + \sqrt{u_Q^2 + 4 D \gamma \eta}}{2 D}, \\
F &= q_L + (D \lambda + u_Q) a_L.
\end{align}

\subsubsection*{Step 5: use a scalar root solve in $j_B$}
The outer solver first scans a grid in $\log j_B$ across the allowed interval
and evaluates $F$.

\begin{itemize}[leftmargin=1.5em]
\item If a sign-changing bracket is found, the code uses
      \texttt{scipy.optimize.root\_scalar(..., method="brentq")}.
\item If no bracket is found, the code falls back to a bounded scalar
      \texttt{least\_squares} fit in $\log j_B$.
\end{itemize}

This gives a robust scalar shooting strategy without requiring the user to
manually pre-bracket the solution in every case.

\subsection{Diagnostics and caching}
The 1D solver contains:
\begin{itemize}[leftmargin=1.5em]
\item a \texttt{trial\_cache} keyed by $(\log j_B, \texttt{with\_profile})$,
\item best-trial tracking based on the smallest $|F|$ found so far,
\item the same \texttt{verb=False / True / "full"} diagnostic convention used by
the 2D solver.
\end{itemize}

\subsection{Returned result dictionary}
The 1D solver returns a dictionary containing at least:
\begin{center}
\begin{tabular}{>{\ttfamily}p{0.27\textwidth}p{0.60\textwidth}}
\toprule
key & meaning \\
\midrule
success & whether the scalar shooting converged \\
message & textual convergence/failure summary \\
jB & solved baryon flux \\
aQstar & fixed interface composition input \\
tail\_residual & scalar asymptotic-tail residual \\
a\_end, q\_end & end-state IVP residual data at $x_{\max}$ \\
lambda & linearized stable-tail decay rate \\
u\_Q & far-right equilibrium velocity $j_B/n_B^Q$ \\
muB\_Q, nB\_Q, nK\_Q & downstream equilibrium state \\
muB\_Qstar, muK\_Qstar, nB\_Qstar & interface state \\
D, eta, gamma, tau & frozen microphysics coefficients \\
branch\_label & fixed compatibility label \texttt{"muK-rich"} for the
positive-$\mu_K$ physical branch \\
\bottomrule
\end{tabular}
\end{center}

If \texttt{return\_profile=True}, the profile arrays
\[
x,\ a,\ q,\ u,\ n_B,\ \mu_B,\ \mu_K
\]
are also included.

\section{The Curve-Extraction Helper}
\subsection{Function}
\texttt{extract\_jB\_curve\_vs\_aQstar(...)}

\subsection{Purpose}
This helper automates the repeated use of the new 1D solver over a list or array
of fixed $a_{Q^\star}$ values. It is designed to build the numerical curve
\[
j_B(a_{Q^\star})
\]
efficiently.

\subsection{Continuation strategy}
The implementation uses simple one-step continuation:
\begin{itemize}[leftmargin=1.5em]
\item the user supplies a sequence \texttt{aQstar\_values},
\item the code solves the first point using \texttt{jB\_guess},
\item every successful point updates the next initial guess via
      \[
      j_B^{\rm guess, next} \leftarrow j_B^{\rm previous\ success}.
      \]
\end{itemize}

This is the intended fast path for scans such as
\[
a_{Q^\star} \in [0,1].
\]

\subsection{Returned data}
The helper returns a structured dictionary containing:
\begin{itemize}[leftmargin=1.5em]
\item the full per-point result list,
\item a boolean success array,
\item the full $a_{Q^\star}$ and $j_B$ arrays,
\item the subset of successful points,
\item the tail residual values,
\item the branch label.
\end{itemize}


\section{April 2026 Speed Update for the 1D Solver}

\subsection{Motivation}
The first version of \texttt{solve\_steady\_front\_1d\_aQstar(...)} was still
very expensive in practice. Although the outer shooting problem had been reduced
to one scalar unknown $j_B$, one evaluation of the tail residual still required:
\begin{itemize}[leftmargin=1.5em]
\item constructing the nuclear endpoint $N$,
\item solving the far-right quark endpoint $Q$,
\item solving the interface state $Q^\star$,
\item integrating the quark-side IVP,
\item repeatedly solving the local quark closure problem along the IVP.
\end{itemize}
The dominant cost was therefore not the number of shooting parameters, but the
cost per residual evaluation and the number of residual evaluations used to find
a scalar bracket in $j_B$.

\subsection{Exact massless finite-temperature thermodynamics}
For the production case $m_s=0$, the quark thermodynamics can avoid numerical
quadrature for massless flavors. The helpers \texttt{P\_f}, \texttt{E\_f}, and
\texttt{n\_B} now short-circuit the case $m=0$ and use the exact formulas
\begin{align}
P_f(\mu,T;m=0) &= \frac{\mu^4}{12\pi^2}
              + \frac{\mu^2 T^2}{6}
              + \frac{7\pi^2 T^4}{180},\\
n_B(\mu,T;m=0) &= \frac{\mu^3}{3\pi^2} + \frac{\mu T^2}{3},\\
E_f(\mu,T;m=0) &= 3P_f(\mu,T;m=0).
\end{align}
These formulas use the same convention as the rest of the file: spin degeneracy
is included and quark color is not included at the single-flavor helper level.
Color is still applied in \texttt{\_quark\_density\_uds(...)} and related quark
matter wrappers.

This change removes most finite-temperature quadrature calls for the current
$m_s=0$ steady-front scans. The massive path is kept unchanged for $m\ne0$.

\subsection{Local bracket expansion in the 1D shooting solve}
The original 1D implementation scanned a broad grid of up to roughly forty
values in $\log j_B$ before attempting a scalar root solve. Since each point in
that scan required a full IVP integration, this was unnecessarily expensive.

The 1D solver now uses local bracket expansion around the supplied
\texttt{jB\_guess}. The algorithm is:
\begin{enumerate}[leftmargin=1.5em]
\item evaluate the tail residual at the initial guess;
\item if it already satisfies the tolerance, accept it;
\item otherwise expand outward multiplicatively using a factor of $1.5$;
\item check the pairs $(j_B/1.5^k,j_B)$ and $(j_B,j_B\,1.5^k)$ until a sign
      change is found or the allowed bounds are reached;
\item use \texttt{root\_scalar(..., method="brentq")} once a sign-changing
      bracket is found;
\item fall back to bounded one-dimensional \texttt{least\_squares} only if no
      bracket is found.
\end{enumerate}
The implementation solves in $\log j_B$ internally, so the positivity of $j_B$
is preserved.

The result dictionary records this root-finding path using
\texttt{\_bracket\_found}, \texttt{\_bracket\_evals}, \texttt{\_root\_method},
and \texttt{shooting\_evals}. These fields make it easier to diagnose whether a
run used the fast bracketing path or the fallback least-squares path.

\subsection{Normalized tail residual}
The raw tail residual,
\[
F_{\rm tail} = q(L) + (D\lambda+u_Q)a(L),
\]
can be numerically tiny simply because the absolute scale of the flux variable
$q$ is tiny. Therefore, using an absolute tolerance on $F_{\rm tail}$ can
incorrectly accept the initial guess without meaningful shooting iteration.

The 1D solver now uses a signed normalized residual for root-finding and
convergence:
\[
\widehat F_{\rm tail}
= \frac{F_{\rm tail}}
       {\max\left(|q(L)|,\ |(D\lambda+u_Q)a(L)|,\ {\rm tiny}\right)}.
\]
The raw dimensional residual is still stored as \texttt{tail\_residual}, while
the normalized residual is stored as \texttt{tail\_residual\_norm}. The scalar
root solve and final success test use \texttt{tail\_residual\_norm}.

\subsection{Fast local continuation for quark closure}
The local quark-state closure solve is performed many times during the IVP. The
closure equations are
\begin{align}
\Pi_{QM}(\mu_B,\mu_K; j_B) - \Pi &= 0,\\
\frac{n_K(\mu_B,\mu_K)-n_K^Q}{n_B^Q} - a &= 0.
\end{align}
The previous implementation always used the broader multi-guess local solver.
That was robust but expensive.

The new implementation adds a single-guess helper,
\texttt{\_solve\_quark\_state\_once\_from\_guess(...)}, and updates
\texttt{\_solve\_local\_quark\_state\_from\_a\_and\_Pi(...)} so that it first
tries the previous successful local solution as the next initial guess. If this
nearest-neighbor continuation solve succeeds and gives a positive-density state,
it returns immediately. Only if this fast continuation step fails does the code
fall back to the broader multi-guess solve on the same positive-$\mu_K$
physical branch.

This keeps the existing positive-$\mu_K$ candidate-ranking logic, but avoids
repeatedly solving from several generic guesses at every IVP RHS call.

\subsection{Lightweight diagnostics}
The 1D solver now stores lightweight performance counters in the returned result
dictionary:
\begin{itemize}[leftmargin=1.5em]
\item \texttt{shooting\_evals}: number of scalar tail-residual evaluations,
\item \texttt{ivp\_rhs\_calls}: number of IVP RHS calls in the best trial,
\item \texttt{q\_root\_calls}: number of far-right $Q$ endpoint solves,
\item \texttt{qstar\_root\_calls}: number of $Q^\star$ root solves,
\item \texttt{local\_state\_calls}: number of local quark closure calls,
\item \texttt{local\_root\_calls}: number of local nonlinear root solves,
\item \texttt{local\_fast\_failures}: number of times nearest-neighbor
      continuation failed and the broad fallback was needed,
\item \texttt{profile\_state\_calls}: number of additional local solves used
      during profile reconstruction,
\item \texttt{\_bracket\_found}: whether the local bracket expansion found a
      sign-changing bracket,
\item \texttt{\_bracket\_evals}: number of evaluations used during bracket
      construction.
\end{itemize}
These fields are intended for performance diagnosis and are also exported by the
cluster scan script summary.

\subsection{Cluster scan-script changes}
The 6Apr cluster script \texttt{new\_paper\_calculations/6Apr/steady\_front\_scan.py}
was updated to make scalar curve extraction the default. By default it now runs
with \texttt{return\_profile=False}; this avoids a second profile reconstruction
and many extra local closure solves for every $a_{Q^\star}$ point.

Profiles can still be generated explicitly:
\begin{itemize}[leftmargin=1.5em]
\item use \texttt{--profile} to reconstruct profiles during the main scan;
\item use \texttt{--profile-successful} to first solve scalar points and then
      rerun successful points once with \texttt{return\_profile=True}.
\end{itemize}
The summary file \texttt{trial\_summary.jsonl} now stores the diagnostic counters
listed above, so cluster runs can be inspected without opening every detailed
log file.

\section{Compact-Coordinate 1D Solver}
\subsection{Function}
The additional function
\[
\texttt{solve\_steady\_front\_1d\_aQstar\_rescaled(...)}
\]
was added immediately after \texttt{solve\_steady\_front\_1d\_aQstar(...)}.
It keeps the same physical state construction and the same scalar shooting
unknown $j_B$ at fixed $a_{Q^\star}$, but integrates the quark-side IVP in a
compact coordinate instead of directly in $x$.

\subsection{Motivation}
The finite-domain 1D solver evaluates the stable-tail residual at a chosen
finite $x_{\max}$. For long tails, increasing $x_{\max}$ to hundreds of nm can
still leave the computed profile far from the downstream equilibrium. The
compact-coordinate solver was added to make the numerical domain better aligned
with the asymptotic tail.

\subsection{What is borrowed from the old numerical velocity code}
The old helper functions \texttt{\_vNtoQ\_num(...)} and
\texttt{\_vNtoQ\_num\_rescale(...)} used a coordinate transformation to map a
semi-infinite interval to a finite computational interval. The new solver uses
only that numerical idea.

It does \emph{not} use the old physical assumptions. In particular, the new
solver does not assume a constant quark-side velocity, does not remove baryon
flux conservation, and does not replace the current hydro closure. The coupled
steady-front equations remain the same as in the existing 1D and 2D solvers:
\begin{align}
j_B &= u(x)n_B(x),\\
\Pi &= h(x)u(x)^2 + P(x),\\
\frac{da}{dx} &= \frac{q+u a}{D},\\
\frac{dq}{dx} &= \gamma(a^3+\eta a).
\end{align}
At every local point, the quark thermodynamic state is still solved from the
fixed values of $(a,\Pi,j_B)$.

\subsection{Coordinate map}
The compact coordinate is
\begin{equation}
s = 1 - e^{-\lambda x/\texttt{kappa\_factor}}, \qquad s \in [0,1),
\end{equation}
so that
\begin{equation}
x(s) = -\frac{\texttt{kappa\_factor}}{\lambda}\log(1-s), \qquad
\frac{dx}{ds} = \frac{\texttt{kappa\_factor}}{\lambda(1-s)}.
\end{equation}
The compactification is still set by the same linearized tail convention used
by the finite-domain 1D solver:
\begin{align}
\lambda &=
\frac{-u_Q + \sqrt{u_Q^2 + 4D\gamma\eta}}{2D},\\
u_Q &= \frac{j_B}{n_B^Q},\\
\frac{1}{\lambda_{\rm eff}}
&\equiv \frac{\texttt{kappa\_factor}}{\lambda}.
\end{align}
Here $u_Q$ is evaluated at the far-right equilibrated state $Q$, while
$D,\eta,\gamma$ are frozen at $Q^\star$, exactly as in the other steady-front
solvers.

\subsection{Transformed equations}
Only the independent variable changes. Applying the chain rule gives
\begin{align}
\frac{da}{ds}
&= \frac{q+u a}{D}\frac{\texttt{kappa\_factor}}{\lambda(1-s)},\\
\frac{dq}{ds}
&= \gamma(a^3+\eta a)\frac{\texttt{kappa\_factor}}{\lambda(1-s)}.
\end{align}
This is the equation system implemented by the new compact solver. It keeps
$u(x)=j_B/n_B(x)$ local and variable; it is not the old constant-velocity
approximation.

\subsection{Endpoint handling}
The mathematical endpoint $s=1$ corresponds to $x=\infty$, so the code never
integrates exactly to $s=1$. A direct choice such as
$s_{\rm end}=1-\texttt{tail\_eps}$ can be too aggressive, because the compact
ODE contains $dx/ds=\texttt{kappa\_factor}/[\lambda(1-s)]$ and therefore becomes singular near
$s=1$. Therefore, when \texttt{x\_max=None}, the code uses a tunable finite
number of compact tail lengths:
\[
s_{\rm end}
= \min\left(1-\texttt{tail\_eps},
1-e^{-\texttt{compact\_tail\_lengths}}\right).
\]
The default is \texttt{compact\_tail\_lengths=0.35}. This is an automatic
compact endpoint, not an attempt to integrate exactly to infinity. To push the
endpoint farther out, increase \texttt{compact\_tail\_lengths} or provide an
explicit \texttt{x\_max}.

Equivalently, when \texttt{x\_max=None},
\[
x_{\rm end}
= -\frac{\texttt{kappa\_factor}}{\lambda}\log(1-s_{\rm end})
\simeq \frac{\texttt{kappa\_factor}}{\lambda}\,\texttt{compact\_tail\_lengths},
\]
unless the hard cap from \texttt{tail\_eps} is reached. Therefore
\texttt{compact\_tail\_lengths} is a dimensionless endpoint-depth parameter
measured in units of the derived length scale
$\texttt{kappa\_factor}/\lambda$. For example:
\begin{center}
\begin{tabular}{ccc}
\toprule
\texttt{compact\_tail\_lengths} & $s_{\rm end}$ & $x_{\rm end}$ \\
\midrule
0.35 & $1-e^{-0.35}\simeq 0.295$ & $0.35\,\texttt{kappa\_factor}/\lambda$ \\
1.0  & $1-e^{-1}\simeq 0.632$    & $1.0\,\texttt{kappa\_factor}/\lambda$ \\
2.0  & $1-e^{-2}\simeq 0.865$    & $2.0\,\texttt{kappa\_factor}/\lambda$ \\
3.0  & $1-e^{-3}\simeq 0.950$    & $3.0\,\texttt{kappa\_factor}/\lambda$ \\
\bottomrule
\end{tabular}
\end{center}

The notebook variable \texttt{compact\_tail\_lengths\_rescaled} is just the
notebook-side name for this same solver argument. Increasing it asks the
compact solver to integrate farther into the tail when \texttt{x\_max=None}.
This may improve asymptotic coverage, but it also makes the IVP harder because
$dx/ds$ grows as $s$ approaches 1. If the local quark closure starts failing
near the far tail, reduce \texttt{compact\_tail\_lengths\_rescaled} or use an
explicit finite \texttt{x\_max}.

If \texttt{x\_max} is supplied, the code computes
\[
s(x_{\max}) = 1 - e^{-\lambda x_{\max}/\texttt{kappa\_factor}}
\]
and uses
\[
s_{\rm end} =
\min\left(1-\texttt{tail\_eps},\ s(x_{\max})\right).
\]
The returned dictionary stores both \texttt{s\_end} and the corresponding
physical \texttt{x\_end}. When \texttt{return\_profile=True}, the returned
\texttt{x} array is reconstructed from the compact coordinate using
$x=-(\texttt{kappa\_factor}/\lambda)\log(1-s)$.

\subsection{Tail residual and shooting}
The compact solver uses the same normalized stable-tail residual as the
finite-domain 1D solver, evaluated at the compact endpoint:
\begin{align}
F_{\rm tail}
&= q(s_{\rm end})
 + \left(D\lambda+u_Q\right)a(s_{\rm end}),\\
\widehat F_{\rm tail}
&=
\frac{F_{\rm tail}}
{\max\left(|q(s_{\rm end})|,
|(D\lambda+u_Q)a(s_{\rm end})|,
{\rm tiny}\right)}.
\end{align}
The scalar shooting variable is still $z=\log j_B$, so the solver enforces
$j_B>0$. It uses the same local bracket expansion around \texttt{jB\_guess},
then \texttt{root\_scalar(..., method="brentq")} when a sign-changing bracket
is found, with bounded \texttt{least\_squares} as the fallback.

\subsection{New controls and returned fields}
The compact solver adds the following controls:
\begin{itemize}[leftmargin=1.5em]
\item \texttt{tail\_eps}: how close the compact integration gets to $s=1$;
\item \texttt{kappa\_factor}: multiplier in the compact map
      $s=1-\exp(-\lambda x/\texttt{kappa\_factor})$;
\item \texttt{compact\_tail\_lengths}: automatic endpoint depth used when
      \texttt{x\_max=None};
\item \texttt{x\_max}: optional physical cutoff, still in the existing natural
      unit convention.
\end{itemize}
In addition to the standard 1D result fields, the compact solver returns
\texttt{lambda}, \texttt{s\_end}, \texttt{x\_end\_target},
\texttt{compact\_tail\_lengths}, and \texttt{coordinate}. If a profile is
requested, the returned arrays include \texttt{s} as well as the reconstructed
physical \texttt{x}. For backward compatibility the returned dictionary also
keeps a derived field \texttt{kappa}, equal to
$\texttt{kappa\_factor}/\lambda$, but the solver itself now uses $\lambda$
directly.

\fi


\section{Recommended Solver: Compact-Coordinate 1D BVP}
\subsection{Function and role}
The recommended current solver is
\[
\texttt{solve\_steady\_front\_1d\_aQstar\_rescale\_bvp(...)}.
\]
It is the cleanest implementation of the steady hydro--diffusion--reaction
front because it solves the front profile and the baryon flux $j_B$ together as
a boundary-value problem. The input $a_{Q^\star}$ is treated as the continuation
or scan parameter; $j_B$ is the unknown scalar parameter adjusted by the BVP
solver.

The function is designed as a drop-in solver for other scripts: provide the
physical inputs $(T,n_B^N,B^{1/4})$, the fixed interface composition
$a_{Q^\star}$. The return
value is a dictionary containing the endpoint states, the BVP diagnostics, and,
if requested, the full profile arrays.

\subsection{Inputs that define the physical problem}
The primary inputs are:
\begin{itemize}[leftmargin=1.5em]
\item \texttt{T}: temperature;
\item \texttt{nB\_N}: upstream nuclear baryon density at $0^-$;
\item \texttt{B\_one\_forth}: bag-model parameter $B^{1/4}$;
\item \texttt{aQstar}: prescribed interface composition $a(0^+)$;
\item \texttt{ms}: strange-quark mass, with production scans usually using
      $m_s=0$;
\end{itemize}
Numerical controls include \texttt{tail\_eps}, \texttt{n\_mesh},
\texttt{tol\_bvp}, \texttt{max\_nodes}, \texttt{jB\_guess},
\texttt{jB\_bounds}, and \texttt{return\_profile}.

\subsection{What the solver borrows from the old numerical code}
The old \texttt{\_vNtoQ\_num(...)} family used a BVP solve and a compact
coordinate. The new solver borrows only that numerical architecture. It does
\emph{not} copy the old physical simplifications. In particular, it does not
assume a constant quark-side velocity and it does not remove baryon-flux
conservation.

The current equations remain
\begin{align}
j_B &= u(x)n_B(x),\\
\Pi &= h(x)u(x)^2 + P(x),\\
\frac{da}{dx} &= \frac{q+u a}{D},\\
\frac{dq}{dx} &= \gamma(a^3+\eta a),
\end{align}
with the local velocity always reconstructed from
\begin{equation}
u(x)=\frac{j_B}{n_B(x)}.
\end{equation}

\subsection{Compact coordinate and transformed ODE}
The compact coordinate is
\begin{equation}
s = 1-e^{-\lambda x/\texttt{kappa\_factor}}, \qquad 0\le s < 1,
\end{equation}
so that
\begin{equation}
x(s)=-\frac{\texttt{kappa\_factor}}{\lambda}\log(1-s),\qquad
\frac{dx}{ds}=\frac{\texttt{kappa\_factor}}{\lambda(1-s)}.
\end{equation}
The BVP endpoint is truncated at
\begin{equation}
s_{\rm end}=1-\texttt{tail\_eps}.
\end{equation}
The tail decay rate is derived from the downstream linear tail:
\begin{align}
u_Q &= \frac{j_B}{n_B^Q},\\
\lambda &= \frac{-u_Q+\sqrt{u_Q^2+4D\gamma\eta}}{2D},\\
\frac{1}{\lambda_{\rm eff}} &= \frac{\texttt{kappa\_factor}}{\lambda}.
\end{align}
Here $u_Q$ and $n_B^Q$ are evaluated at the equilibrated far-right quark state
$Q$, while $D$, $\eta$, and $\gamma$ are frozen at the interface state
$Q^\star$.

By the chain rule the equations supplied to \texttt{solve\_bvp} are
\begin{align}
\frac{da}{ds}
&= \frac{q+u a}{D}\frac{\texttt{kappa\_factor}}{\lambda(1-s)},\\
\frac{dq}{ds}
&= \gamma(a^3+\eta a)\frac{\texttt{kappa\_factor}}{\lambda(1-s)}.
\end{align}
This is only a coordinate change; the physics and closure are unchanged.

\subsection{Unknown BVP parameter and bounds}
The unknown BVP parameter controls $j_B$. If no explicit bounds are supplied,
\begin{equation}
j_B=e^\theta,
\end{equation}
which guarantees $j_B>0$. If \texttt{jB\_bounds=(j\_min,j\_max)} is supplied,
the code uses the bounded map
\begin{equation}
j_B=j_{\min}+(j_{\max}-j_{\min})\frac{1}{1+e^{-\theta}}.
\end{equation}
For stable scans, the upper bound should not be set absurdly high. In the 6Apr
runs, allowing \texttt{jB\_max=50} let the BVP escape to an unphysical upper
boundary; the conservative scan default is now \texttt{jB\_bounds=1e-5,5.0}.

\subsection{Endpoint-state construction helper chain}
For each trial value of $j_B$, the solver constructs the global state as follows:
\begin{enumerate}[leftmargin=1.5em]
\item Compute the upstream nuclear state $N$ from $(T,n_B^N)$ using
      \texttt{PNM\_n(...)} and \texttt{edensNM\_n(...)}. Then
      \[
      u_N=\frac{j_B}{n_B^N},\qquad \Pi=h_Nu_N^2+P_N.
      \]
\item Use \texttt{\_solve\_muB\_Q\_at\_muK0\_for\_given\_Pi(...)} to find the
      downstream equilibrated quark state $Q$ at $\mu_K=0$.
\item Compute $n_B^Q$ with \texttt{nB\_QM(...)} and $n_K^Q$ with
      \texttt{nK\_QM(...)}. For $m_s=0$, the code sets $n_K^Q=0$ exactly by
      symmetry at $\mu_K=0$.
\item Compute
      \[
      a_N=\frac{n_B^N-n_K^Q}{n_B^Q}.
      \]
\item Use the interface-state helper
      \codepath{_solve_interface_Qstar_from_aQstar_and_Pi(...)} to find
      $(\mu_B^{Q^\star},\mu_K^{Q^\star})$ from momentum-flux matching and
      the fixed composition condition
      \[
      \frac{n_K(\mu_B,\mu_K)-n_K^Q}{n_B^Q}=a_{Q^\star}.
      \]
\item Use \texttt{\_microphysics\_at\_Qstar(...)} to compute the frozen
      coefficients $D$, $\eta$, $\gamma$, and $\tau$.
\end{enumerate}
This helper chain is the reusable part of the implementation: any external
copy of the method should keep this state-construction order unchanged.

\subsection{Local hydro closure inside the BVP}
At each BVP mesh point, the local profile value $a(s)$ is known but the quark
thermodynamic state is not. The solver calls
\texttt{\_solve\_local\_quark\_state\_from\_a\_and\_Pi(...)} to solve
\begin{align}
\Pi_{QM}(\mu_B,\mu_K;j_B)-\Pi &= 0,\\
\frac{n_K(\mu_B,\mu_K)-n_K^Q}{n_B^Q}-a(s) &= 0.
\end{align}
Then it computes
\begin{equation}
n_B(s)=n_B^{QM}(\mu_B,\mu_K),\qquad
u(s)=\frac{j_B}{n_B(s)}.
\end{equation}
The previous mesh-point solution is used as the next initial guess whenever
possible, preserving nearest-neighbor continuation and reducing accidental
branch jumps.

\subsection{Boundary conditions}
Because the system has two first-order ODE variables and one unknown scalar
parameter, the BVP needs three boundary conditions:
\begin{align}
a(0)-a_{Q^\star} &= 0,\\
q(0)+a_Nu_N &= 0,\\
q(s_{\rm end})+(D\lambda+u_Q)a(s_{\rm end}) &= 0.
\end{align}
The first fixes the input interface composition. The second is the interface
flux jump condition
\[
q(0^+)=-a_Nu_N.
\]
The third enforces the downstream stable-tail manifold.

\subsection{Physical branch convention}
The current steady-front implementation does not expose a
\texttt{muK\_rich=True/False} switch or a deliberate $\mu_K$-poor mode. The
helper \texttt{\_branch\_muK\_seed(...)} supplies only positive-$\mu_K$ seeds,
and the interface/local closure helpers rank candidate roots against those
positive seeds. In other words, the code follows a single positive-$\mu_K$
physical branch. The returned string \texttt{branch\_label="muK-rich"} is
retained only as a compatibility label in result dictionaries.

\subsection{Returned dictionary}
The BVP solver returns a dictionary containing, at minimum,
\texttt{success}, \texttt{message}, \texttt{jB}, \texttt{aQstar},
\texttt{a\_N}, \texttt{u\_N}, \texttt{u\_Q}, \texttt{Pi},
\texttt{muB\_Q}, \texttt{nB\_Q}, \texttt{nK\_Q}, \texttt{muB\_Qstar},
\texttt{muK\_Qstar}, \texttt{nB\_Qstar}, \texttt{D}, \texttt{eta},
\texttt{gamma}, \texttt{tau}, \texttt{lambda},
\texttt{s\_end}, \texttt{x\_end}, \texttt{tail\_residual},
\texttt{tail\_residual\_norm}, and BVP diagnostics such as
\texttt{bvp\_status}, \texttt{bvp\_message}, \texttt{bvp\_niter}, and
\texttt{bvp\_nodes}. For backward compatibility the result also keeps a
derived field \texttt{kappa}, but internally the compact map is written
entirely in terms of $\lambda$.

If \texttt{return\_profile=True}, it also returns arrays
\texttt{s}, \texttt{x}, \texttt{a}, \texttt{q}, \texttt{u}, \texttt{nB},
\texttt{muB}, and \texttt{muK}.

\subsection{Recommended scan usage}
The 6Apr scan driver should follow the branch by continuation rather than
starting each $a_{Q^\star}$ in an isolated worker. The conservative recommended
command is
\begin{verbatim}
python steady_front_scan.py \
  --workers 1 \
  --chunk-size 20 \
  --jB-bounds 1e-5,5.0 \
  --jB-aQstar-slope 1.2 \
  --tail-eps 1e-8 \
  --verb-mode simple
\end{verbatim}
The scan script now solves the scalar BVP first and only reconstructs profiles
after scalar success. This prevents a late profile-reconstruction failure from
hiding an otherwise useful scalar result for $j_B$.

\clearpage
\section{Entropy-Enabled Compact BVP Solver}
\subsection{Purpose and relation to the older compact BVP}
The function
\texttt{solve\_steady\_front\_1d\_aQstar\_\allowbreak rescale\_bvp(...)}
remains the non-entropy compact-coordinate BVP path. It solves for
$(a,q)$ together with the unknown scalar flux $j_B$, uses the fixed
temperature input $T$, and freezes $D$, $\eta$, and $\gamma$ at the
interface state $Q^\star$.

The newer function \texttt{solve\_steady\_front\_entropy(...)} keeps the same
overall compact-coordinate architecture but enlarges the state to include
entropy production and a variable temperature profile. The ODE variables are
\begin{equation}
a(x),\qquad q(x),\qquad w(x)=s(x)u(x),
\end{equation}
and the local thermodynamic closure is upgraded from
$(\mu_B,\mu_K)$ to $(\mu_B,\mu_K,\log T)$. In this entropy-enabled path the
temperature is reconstructed pointwise, and the microphysics coefficients are
evaluated locally rather than frozen once at $Q^\star$.

The additional function \texttt{solve\_steady\_front\_entropy\_TQguess(...)}
keeps the same three ODE fields but upgrades the global BVP solve from one
unknown scalar parameter to two: $(j_B,T_Q)$. It still treats the interface
temperature as the input $T$, but it no longer forces the far-right
equilibrated quark state $Q$ to use that same temperature.

\subsection{Helper chain and solver workflow}
The entropy-enabled implementation is organized as a reusable helper chain:
\begin{enumerate}[leftmargin=1.5em]
\item \texttt{\_quark\_thermo\_state(...)} builds a fully
      \texttt{ms}-consistent local quark state from
      $(\mu_B,\mu_K,T,j_B)$ and returns
      $n_B$, $n_K$, $P$, $e$, $h$, $u$, $s$, $w=s\,u$, and $\Pi$.
\item \texttt{\_solve\_muB\_Q\_at\_muK0\_for\_given\_Pi\_ms(...)} solves for
      the equilibrated far-right quark state $Q$ at $\mu_K=0$ using the
      \texttt{ms}-aware momentum-flux helper throughout.
\item \texttt{\_microphysics\_from\_quark\_state(...)} computes the local
      coefficients $D$, $\eta$, $\gamma$, and $\tau$ from the local
      thermodynamic state, with $\mu_Q=\mu_B/3$ and the local temperature.
\item \texttt{\_quark\_state\_entropy\_residual(...)} and
      \texttt{\_quark\_state\_entropy\_residual\_ok(...)} define the
      three-equation entropy closure residual and the acceptance test for a
      local root solve.
\item \texttt{\_solve\_quark\_entropy\_state\_\allowbreak once\_from\_guess(...)}
      carries out one nonlinear solve in $(\mu_B,\mu_K,\log T)$ from a single
      continuation guess.
\item \texttt{\_solve\_local\_quark\_state\_from\_a\_w\_and\_Pi(...)} wraps
      the previous helper into the local closure used inside the BVP. It
      first tries nearest-neighbor continuation, then falls back to a richer
      seed set on the same physical positive-$\mu_K$ branch.
\item \texttt{\_seed\_entropy\_jB\_guess(...)} obtains an initial guess
      for $j_B$ by reusing the current non-entropy compact BVP solver on a
      small-$a_{Q^\star}$ seed problem.
\item \texttt{\_solve\_steady\_front\_entropy\_once(...)} constructs the
      endpoint states, solves the direct entropy-enabled BVP once, and
      optionally reconstructs the profile arrays.
\item \texttt{solve\_steady\_front\_entropy(...)} is the public wrapper. It
      first tries the direct one-shot solve, and if that fails, it switches to
      staged continuation in $a_{Q^\star}$ while recycling the last converged
      entropy profile as the next initial guess.
\item \texttt{\_solve\_steady\_front\_entropy\_TQguess\_once(...)} is the
      direct two-parameter entropy BVP with global unknowns $(j_B,T_Q)$.
\item \texttt{solve\_steady\_front\_entropy\_TQguess(...)} is the public
      wrapper for that two-parameter problem. It mirrors the continuation
      logic of \texttt{solve\_steady\_front\_entropy(...)} but carries both
      $j_B$ and $T_Q$ forward between stages.
\end{enumerate}
The small utility \texttt{\_strip\_entropy\_profile\_fields(...)} only trims
profile arrays from returned dictionaries when the wrapper does not need them;
it is a return-shaping helper, not a physics helper.

\subsection{Thermodynamic closure and local state reconstruction}
At fixed $j_B$, $\Pi$, $a$, and $w$, the local entropy-enabled quark state is
reconstructed by solving for $(\mu_B,\mu_K,\log T)$ rather than
$(\mu_B,\mu_K,T)$ directly. The closure equations are
\begin{align}
\Pi_{QM}(\mu_B,\mu_K,T;j_B)-\Pi &= 0,\\
\frac{n_K(\mu_B,\mu_K,T)-n_K^Q}{n_B^Q}-a &= 0,\\
s_{QM}(\mu_B,\mu_K,T)\,\frac{j_B}{n_B^{QM}(\mu_B,\mu_K,T)}-w &= 0.
\end{align}
Here $n_B^Q$ and $n_K^Q$ come from the asymptotic equilibrated state $Q$.

The entropy density is not computed by a separate EOS entry point. Instead the
code reconstructs it from the thermal relation
\begin{equation}
s=\frac{e+P-\mu_B n_B-\mu_K n_K}{T},
\end{equation}
using the existing quark EOS helpers for $e$, $P$, $n_B$, and $n_K$.
This keeps the entropy-enabled path fully \texttt{ms}-consistent.

The entropy solver fails fast on nonphysical local states. In particular, the
new helper chain rejects $T\le 0$, $w\le 0$, non-positive $n_B$, negative
$\mu_K$ beyond numerical tolerance, non-positive entropy density, and local
root solves that do not satisfy the momentum, composition, and entropy-flux
constraints to the required tolerance.

\subsection{Local microphysics and compact-coordinate BVP}
The compact coordinate now uses the code name \texttt{s\_coord} to avoid
confusion with the entropy density $s$:
\begin{equation}
s_{\rm coord}=1-e^{-\lambda x},\qquad
x(s_{\rm coord})=-\frac{1}{\lambda}\log(1-s_{\rm coord}),\qquad
\frac{dx}{ds_{\rm coord}}=\frac{1}{\lambda(1-s_{\rm coord})}.
\end{equation}
The entropy-enabled profile evolves the first-order system
\begin{align}
\frac{da}{dx} &= \frac{q+u a}{D(x)},\\
\frac{dq}{dx} &= \gamma(x)\bigl(a^3+\eta(x)a\bigr),\\
\frac{dw}{dx} &= \frac{n_B^Q}{T(x)}\mu_K(x)\gamma(x)\bigl(a^3+\eta(x)a\bigr),
\end{align}
with $w=s\,u$ and
\begin{equation}
u(x)=\frac{j_B}{n_B(x)}.
\end{equation}
After the compact-coordinate change this becomes
\begin{align}
\frac{da}{ds_{\rm coord}}
&= \frac{q+u a}{D(x)}\frac{1}{\lambda(1-s_{\rm coord})},\\
\frac{dq}{ds_{\rm coord}}
&= \gamma(x)\bigl(a^3+\eta(x)a\bigr)\frac{1}{\lambda(1-s_{\rm coord})},\\
\frac{dw}{ds_{\rm coord}}
&= \frac{n_B^Q}{T(x)}\mu_K(x)\gamma(x)\bigl(a^3+\eta(x)a\bigr)
\frac{1}{\lambda(1-s_{\rm coord})}.
\end{align}
Unlike the older compact BVP, the coefficients
\begin{equation}
D(x),\qquad \eta(x),\qquad \gamma(x),\qquad \tau(x)
\end{equation}
are evaluated pointwise from the local state through
\texttt{\_microphysics\_from\_quark\_state(...)}.

The original entropy-enabled BVP has three first-order fields and one unknown
scalar parameter, so it uses four boundary conditions:
\begin{align}
a(0)-a_{Q^\star} &= 0,\\
q(0)+a_Nu_N &= 0,\\
w(0)-s_{Q^\star}u_{Q^\star} &= 0,\\
q(s_{\rm end})+\bigl(D_Q\lambda+u_Q\bigr)a(s_{\rm end}) &= 0.
\end{align}
The first two conditions are the same composition and interface-flux
constraints as before. The third enforces interface entropy-flux continuity on
the quark side, using the interface state $Q^\star$ at the input temperature.
The fourth is the downstream stable-tail condition. In the current
implementation the tail quantities are built from the asymptotic equilibrated
state $Q$:
\begin{align}
u_Q &= \frac{j_B}{n_B^Q},\\
\lambda &= \frac{-u_Q+\sqrt{u_Q^2+4D_Q\gamma_Q\eta_Q}}{2D_Q},
\end{align}
where $D_Q$, $\eta_Q$, and $\gamma_Q$ are evaluated at $Q$, not at
$Q^\star$.

\subsection{What is new in \texttt{solve\_steady\_front\_entropy\_TQguess(...)}}
The newer two-parameter solver keeps the same ODE fields $(a,q,w)$ and the
same local closure in $(\mu_B,\mu_K,\log T)$, but it promotes the downstream
equilibrium temperature $T_Q$ to a second global BVP unknown. Its public
interface is
\begin{verbatim}
solve_steady_front_entropy_TQguess(
    T, nB_N, B_one_forth, aQstar,
    ms=0.0, param=para.paraQMCRMF3, NM_type="PNM",
    tail_eps=1e-8, n_mesh=200, tol_bvp=1e-4, max_nodes=10000,
    jB_guess=None, TQ_guess=None, jB_bounds=None,
    return_profile=False, verb=False,
)
\end{verbatim}
The input temperature $T$ still defines the upstream state $N(0^-)$ and the
interface quark state $Q^\star(0^+)$, but the asymptotic state $Q(+\infty)$
is rebuilt from the current trial pair $(j_B,T_Q)$:
\begin{enumerate}[leftmargin=1.5em]
\item compute the conserved momentum flux $\Pi$ from the upstream state,
\item impose $\mu_K^Q=0$,
\item solve the scalar momentum-flux equation for $\mu_B^Q$ at the trial $T_Q$,
\item reconstruct $n_B^Q$, $n_K^Q$, $u_Q$, $s_Q$, and
      $w_Q=s_Q u_Q$ from that equilibrium state,
\item rebuild the asymptotic microphysics and tail coefficient from that same
      corrected $Q$ state.
\end{enumerate}

Because the problem now has three first-order fields and two unknown global
parameters $(j_B,T_Q)$, the BVP uses five boundary conditions:
\begin{align}
a(0)-a_{Q^\star} &= 0,\\
q(0)+a_Nu_N &= 0,\\
w(0)-s_{Q^\star}u_{Q^\star} &= 0,\\
q(s_{\rm end})+\bigl(D_Q\lambda+u_Q\bigr)a(s_{\rm end}) &= 0,\\
w(s_{\rm end})-s_Qu_Q &= 0.
\end{align}
The last condition is the new entropy-tail matching equation that determines
the solved downstream temperature together with the profile.

The returned dictionary keeps the same scalar and profile fields as
\texttt{solve\_steady\_front\_entropy(...)} and adds
\texttt{T\_Q}, \texttt{w\_Q}, \texttt{s\_Q},
\texttt{entropy\_tail\_residual}, and
\texttt{entropy\_tail\_residual\_norm}. When continuation is used, the wrapper
now carries forward both the last converged $j_B$ and the last converged
$T_Q$, together with the previous $(a,q,w)$ profile.

\subsection{One-shot solver, public wrapper, and returned data}
The public entropy solver has the interface
\begin{verbatim}
solve_steady_front_entropy(
    T, nB_N, B_one_forth, aQstar,
    ms=0.0, param=para.paraQMCRMF3, NM_type="PNM",
    tail_eps=1e-8, n_mesh=200, tol_bvp=1e-4, max_nodes=10000,
    jB_guess=None, jB_bounds=None, return_profile=False, verb=False,
)
\end{verbatim}
Its leading arguments mirror the older compact BVP path: upstream
thermodynamic input, bag parameter, and the fixed interface composition
$a_{Q^\star}$. If \texttt{return\_profile=True}, the returned dictionary also
includes the profile arrays
\texttt{s\_coord}, \texttt{x}, \texttt{a}, \texttt{q}, \texttt{w},
\texttt{u}, \texttt{nB}, \texttt{muB}, \texttt{muK},
\texttt{T\_profile}, \texttt{s\_profile},
\texttt{D\_profile}, \texttt{eta\_profile}, and
\texttt{gamma\_profile}.

At minimum, the returned scalar diagnostics include the endpoint-state fields
\texttt{success}, \texttt{message}, \texttt{jB}, \texttt{aQstar},
\texttt{a\_N}, \texttt{u\_N}, \texttt{u\_Q}, \texttt{Pi},
\texttt{muB\_Q}, \texttt{nB\_Q}, \texttt{nK\_Q}, \texttt{muB\_Qstar},
\texttt{muK\_Qstar}, \texttt{nB\_Qstar}, and \texttt{s\_Qstar}; the
asymptotic-tail and microphysics fields \texttt{D\_Q}, \texttt{eta\_Q},
\texttt{gamma\_Q}, \texttt{tau\_Q}, \texttt{lambda},
\texttt{s\_end}, \texttt{x\_end}, \texttt{tail\_residual}, and
\texttt{tail\_residual\_norm}; and the BVP diagnostics
\texttt{bvp\_status}, \texttt{bvp\_message}, \texttt{bvp\_niter}, and
\texttt{bvp\_nodes}. When continuation is used, the wrapper also returns
\texttt{continuation\_used}, \texttt{continuation\_steps},
\texttt{continuation\_seed\_aQstar}, and
\texttt{continuation\_refined}. The result dictionaries still include a
derived compatibility field \texttt{kappa=1/lambda}, but the BVP itself is
written directly in terms of $\lambda$.

The helper \texttt{\_solve\_steady\_front\_entropy\_once(...)} is the direct
one-shot BVP solver. The public wrapper
\texttt{solve\_steady\_front\_entropy(...)} first tries that direct path. If
it fails, the wrapper seeds $j_B$ from a small-$a_{Q^\star}$ non-entropy BVP,
solves a sequence of coarse continuation stages in $a_{Q^\star}$ while
recycling the last converged profile, and finally retries the full target
solve with the requested mesh and tolerance. The two-parameter public wrapper
\texttt{solve\_steady\_front\_entropy\_TQguess(...)} uses the same pattern,
but carries both $(j_B,T_Q)$ between stages. All of these entropy-enabled
solvers follow only the physical positive-$\mu_K$
branch. The compatibility field \texttt{branch\_label} is still reported as
\texttt{"muK-rich"}.

\paragraph{Current limits.}
The continuation and seeding logic substantially improve robustness, but the
current entropy-enabled solver should still be treated as a first
implementation. For sufficiently large $a_{Q^\star}$ the continuation chain
can still fail because the local thermodynamic closure or the BVP collocation
step leaves the physical positive-$\mu_K$ branch. The manual therefore documents
the present behavior as implemented, not a guarantee of universal convergence.

\section{Practical Notes}
\subsection{Physical branch convention}
The current steady-front solvers in \texttt{phase\_velocity.py} do not expose a
separate branch switch. Instead they seed and rank roots so that the solution
stays on the physical positive-$\mu_K$ branch. Where result dictionaries still
carry \texttt{branch\_label}, the value \texttt{"muK-rich"} should be read as a
legacy compatibility label rather than a live branch-selection control.

\subsection{\texorpdfstring{$m_s = 0$ convention at the far-right state}{ms = 0 convention at the far-right state}}
At the downstream equilibrated state $Q$, the code now sets
\[
n_K^Q = 0
\]
exactly whenever $|m_s| \le 10^{-12}$. This is consistent with the intended
symmetry of the massless-flavor limit at $\mu_K = 0$.

\subsection{Recommended scan usage}
If the user wants a one-dimensional solution family over the interval
$a_{Q^\star} \in [0,1]$, the intended pattern is:
\begin{equation}
\texttt{aQstar\_values = np.linspace(0.0, 1.0, N)}.
\end{equation}
The curve-extraction helper will then carry out continuation in $j_B$.

\section{Summary}
The update to \texttt{phase\_velocity.py} introduced steady-front solvers built
around the asymptotic stable-tail condition, including both one-parameter and
two-parameter compact-coordinate BVP paths.
The current implementation:
\begin{itemize}[leftmargin=1.5em]
\item keeps the existing 2D solver intact,
\item reuses the same endpoint-state and local-state construction logic,
\item uses positive-$\mu_K$ seeding and candidate ranking to stay on the
      physical branch,
\item solves a scalar equation for $j_B$ at fixed $a_{Q^\star}$,
\item adds a compact-coordinate 1D variant for long-tail integrations,
\item adds a non-entropy compact-coordinate BVP variant that solves the
      profile and $j_B$ together with \texttt{solve\_bvp},
\item adds an entropy-enabled compact-coordinate BVP variant with local
      closure in $(\mu_B,\mu_K,\log T)$ and entropy flux
      $w=s\,u$ as a third ODE field,
\item updates the quark-side microphysics to be evaluated locally along the
      entropy-enabled profile rather than frozen once at $Q^\star$,
\item adds a two-parameter entropy-enabled compact-coordinate BVP variant
      \texttt{solve\_steady\_front\_entropy\_TQguess(...)} that solves for the
      downstream equilibrium temperature $T_Q$ together with $j_B$,
\item writes the compact-coordinate algebra directly in terms of the tail
      decay rate $\lambda$, while keeping \texttt{kappa} only as a derived
      compatibility field in returned dictionaries,
\item wraps the public entropy solver in seeding and continuation logic so it
      can retry difficult targets with recycled profile guesses,
\item provides a dedicated helper to extract the numerical curve
      $j_B(a_{Q^\star})$ efficiently.
\end{itemize}

\end{document}
